\chapter{Introduction}
\section{Project concept}
This project implements \emph{CM3020 Artificial Intelligence, Project Idea 2: Financial Advisor Bot}. The template requires analysis of financial data, dynamic investment recommendations, a non-technical interface, explanations and evidence of evaluation. The chosen domain is the stock market. The system is educational decision support and simulated portfolio analysis; it does not place orders or provide regulated personal financial advice.

The problem is that a retail investor must combine noisy time-series data, forecasts, portfolio constraints and risk judgements. A chart does not explain how a portfolio should be formed, while an opaque recommendation is difficult to trust. The project treats advice as a pipeline: validated data, forecast, constrained allocation, chronological testing and a plain-language explanation identifying dates and limitations.

\section{Users and objectives}
The target user is a non-specialist investor exploring a small stock universe. The Shiny application shows market data, summary information, allocations, forecasts and evaluation metrics. Objectives were to build a reproducible data path; evaluate reservoir computing against transparent baselines; implement constrained policies including FinRL A2C; create a testable \texttt{AdvisorService}; provide grounded explanations; and evaluate behaviour without look-ahead leakage.

\section{Scope}
The final product is an offline-reproducible MVP using daily data, simulated weights and historical backtesting. Live order execution, intraday data, automatic retraining, user accounts and a guaranteed 2025 test period are outside scope. Production code is under \texttt{Code/advisor}, with \texttt{Code/stock-app/app-express.py} as the UI entry point.
\section{Alignment with the assignment brief}
The template defines this project as CM3020 Artificial Intelligence, Project Idea 2: Financial Advisor Bot. Its central requirement is an integrated and tested system that analyses financial data, produces dynamic investment recommendations, explains those recommendations to a non-technical user and presents evidence that the advice has been evaluated. The final system meets this structure through five connected layers: data ingestion, forecasting, portfolio policy, explanation and web presentation.

The preliminary report proposed Yahoo Finance, Bloomberg and CCXT as possible sources, reservoir computing for prediction, FinRL for allocation, Qwen3 and Smolagents for language interaction, and backtesting for evidence. The final project narrows this proposal into a reproducible stock-market MVP. Yahoo/cache and canonical CSV data are supported; the ESN is evaluated against baselines; an A2C policy is trained and saved for a labelled legacy period; Qwen is restricted to validated facts; and the complete workflow runs through Python Shiny. This refinement responds to the brief's distinction between a model-training system and the interface used by a non-technical user.

The project is therefore assessed as a complete educational decision-support system rather than as a claim of guaranteed investment performance. It demonstrates software engineering, data science, machine learning, web programming and evaluation in one coherent workflow.
